\section{Restricting Speculative Leakage}
\label{sec:controlledspec}
With attacker observations restricted to \sw{shm}, we can now focus on enforcing the execution model.
Observable traces when running a given program on our design should not leak more information than when running the same program non-speculatively i.e. using \sw{seq}.
This can be achieved through a combination of mechanisms for controlled speculation and program analysis.
We show two different approaches: preventing speculative shared-memory accesses altogether or ensuring that their addresses do not leak secrets.
We illustrate these approaches by building two minimal hardware mechanisms to enforce RMI for our enclave programs.
Safe mode achieves RMI without requiring any program analysis and Burst mode enables better performance on specific code snippets at the cost or requiring minimal code analysis.

\begin{figure}[t]
\centering
\includegraphics[width=\columnwidth]{figures/MemExePipeline.pdf}
\caption{Changes to the Memory Execution Pipeline in Citadel to enable Safe mode. New elements are shown in blue.}
\label{fig:MemExe}
\centering
\includegraphics[width=\columnwidth]{figures/FetchPipeline.pdf}
\caption{Changes to the Fetch Pipeline in Citadel to enable Burst mode. New elements are shown in blue.}
\label{fig:fetchpipe}
\end{figure}

\subsection{Safe mode: Preventing Speculative Shared Memory}
\label{sec:mechanism1}
Our first approach consists of preventing speculative shared memory accesses altogether, effectively removing any speculative transmitters that can reach the attacker.
This approach enforces RMI for any program without extra program analysis, making it a great foundational mode to use by default.
We illustrate this principle by building a simple mechanism.
\emph{Safe} mode identifies unsafe shared memory accesses early in the memory execution pipeline, delaying their translation and processing until they become non-speculative (i.e., reach the head of the reorder buffer (ROB)).
This is similar \emph{in spirit} to delay-on-miss~\cite{sakalis2019efficient} but only delays \emph{shared} memory accesses and cannot leverage the L1 cache for reasons expressed in section \ref{sec:disentangle}.
A primer on out-of-order processors can be found in Appendix \ref{apendixooo} and microarchitectural details for Safe mode implementation can be found in Appendix \ref{app:safemode}.
Because the private memory range is defined in virtual memory, we can perform this check without any address translation, preventing information leakage through shared microarchitecture.
Safe instructions proceed to the next \MemExePipeline{} stage; unsafe ones are squashed and signaled to the reservation station so they are re-dispatched once they become non-speculative.
\\

\noindent\textbf{Security Analysis}
An attacker is only able to observe the trace of non-speculative accesses to shared memory. 
As a result, under our threat model, $\uarch[safe]{\cdot}{} \vdash \contract[shm]{\cdot}{seq}$ i.e. \textbf{Safe mode satisfies RMI}.

\subsection{Burst mode: Eliminating Information Leakage}
\label{sec:mechanism2}
Our second approach, called \emph{Burst} mode, is a hardware-software co-design that can improve performance for specific code snippets.
It allows some speculative accesses to shared memory (i.e. $\uarch[burst]{\cdot}{} \vdash \contract[shm]{\cdot}{stl}$) but ensures that it only runs programs that will not leak secrets (i.e. only runs $p$ if $p \vdash NI(\contract[shm]{\cdot}{seq} \Rightarrow \contract[shm]{\cdot}{stl})$).
This last step requires some static code analysis, a sometimes difficult problem (see Section \ref{sec:relatedwork}).
In order to simplify this analysis, we apply two simple design principles: 1) we ensure the (speculative) control flow of our code snippet is self-contained 2) we constrain speculation to a simple model where it is not attacker-controlled. 
The first principle indicates that we should deactivate any predictors that enable speculative jumps to arbitrary locations and guard our code snippets with speculative barriers.
The second stems from the fact that disabling more predictors helps simplify program analysis.
Following these principles, we implement a simple static branch predictor, which enables shared memory pipelining on Citadel (disabled in \emph{Safe mode}).
According to our first guidelines, we disable the branch target buffer (BTB), and the return address stack (RAS) (Figure \ref{fig:fetchpipe}).
This is done on demand by writing to a new CSR, using a write instruction that also doubles as a speculation barrier.
Following our second principle, we also disable the branch history table (BHT), effectively implementing straight-line speculation \cite{straightlinespecARM, guarnieri2021hardware} where the only prediction is the hardcoded ``\texttt{pc+4}'', predicting the next instruction without a change in control flow.

\subsection{Code Analysis for Burst mode}
Not all programs satisfy RMI when run in Burst mode. Take the following pathological code:
\noindent\begin{minipage}[]{0.25\textwidth}
\begin{lstlisting}[language={[RISC-V]Assembler}]
  jal far_away_function
  lw a0, 0(secret_address)
\end{lstlisting}
\end{minipage}
Here the secret address might be leaked under Burst mode while it would never be accessed during non-speculative execution.
To use Burst mode safely, we must perform code analysis to ensure snippets are: 1) self-contained to bound the analysis and 2) leak no secrets under straight-line speculation.
Self-containment is verified by checking that control flow remains confined between the instructions turning Burst mode on and off, i.e. no indirect branch and all branch destinations are valid and within the code snippet.
This constraint simplifies our analysis by limiting its scope to the code snippet.
We then verify that the code does not leak secrets under straight-line speculation.
Load, store, and branch instructions are considered potential speculative transmitters and the accessed register are marked as leaked.
For each transmitter, we perform a backward pass to determine dependencies of the leaked value and marking them as leaked.
This pass explores all possible earlier execution paths, including speculative and non-speculative paths when control flow deviates from pc+4.
Our simple speculation model prevents an explosion in possibilities, making this analysis tractable.
If a leaked value depends on another memory value, verification conservatively fails to prevent Spectre-like gadgets.
On each path, dependencies are \emph{declassified} as public if their values are leaked through a non-speculative load or store (we conservatively assume no declassification through branches).
We also prune paths that have been already explored if the set of leaked registers has not changed, and execution is not speculative, as we are not at risk of missing new leakage.
The backward pass yields a set of registers whose initial values could potentially be leaked, along with the corresponding speculative execution paths.
We can then identify conditions under which these paths are incorrect (i.e., not followed in non-speculative execution).
A code snippet fails verification if any non-public value can be leaked.
We developed a Python tool to automate this analysis.

\begin{figure}[htbp]
\begin{minipage}[]{0.25\textwidth}
\begin{lstlisting}[language={[RISC-V]Assembler},
  numbers=left,
  numberstyle=\tiny,
  numbersep=5pt]
  csrwi MSPEC, BURST_ON
  add	a2,a0,a2
  bgeu	a0,a2,.end # len!=0
.loop:
  lbu	a4,0(a1) # leaks src
  add	a1,a1,1
  add	a0,a0,1
  sb	a4,-1(a0) # leaks dest
  bne	a1,a2,.loop #leaks len
.end:
  csrwi MSPEC, BURST_OFF
\end{lstlisting}
\end{minipage}
\begin{minipage}[]{0.22\textwidth}
\begin{lstlisting}[language={[RISC-V]Assembler},
  numbers=left,
  numberstyle=\tiny,
  numbersep=5pt]
  add	a2,a0,a2
  bgeu	a0,a2,.end # len!=0
  csrwi MSPEC, BURST_ON
.loop:
  lbu	a4,0(a1)
  add	a1,a1,1
  add	a0,a0,1
  sb	a4,-1(a0)
  bne	a1,a2,.loop
  csrwi MSPEC, BURST_OFF
.end:
\end{lstlisting}
\end{minipage}
\caption{RISC-V Assembly code for 
\code{memcpy(a0=dest,a1=src,a2=len)}. Left hand-side does not satisfy RMI under Burst mode while the right-hand side does.}
\label{fig:memcpy-assembly}
\end{figure}

\noindent\textbf{\texttt{memcpy} example}
The default assembly can be seen on the left-hand side of Figure \ref{fig:memcpy-assembly}.
Here our tool identifies 3 potential transmitters line 5, 8 and 9.
Let's perform the backward pass for the branch on line 9 initially marking registers \code{a1} and \code{a2} as leaked.
We roll-back execution through line 8 to 6.
Line 6 overwrites \code{a1} which is declassified but replaced by its dependencies i.e. \code{a1}.
On line 5, a transmitter accesses \code{a1} but non-speculatively so it is not declassified.
At line 5, we encounter two possible paths: one from line 9 (which we prune as it is already explored), and another speculative path from line 3.
The analysis reveals a potential straight-line execution path from line 1 to 9 where the initial values of \code{a1} and \code{a2} (i.e., \code{src} and \code{len}) could be leaked.
This vulnerability occurs as Burst mode incorrectly speculates past the branch on line 3, which happens when \code{len==0}.
Our tool also reveals other leaks for \code{dest} and \code{src} under the same conditions.
We add a guard before entering Burst mode to check \code{len!=0} which generates safe assembly (right hand side).
The tool also detected self-containment issues with over-optimized compiler output (loop unrolling and rogue code block placements) where Burst mode would have been insecure.

\noindent\textbf{More Complex Analysis Tools}
The conditions we check on our code snippets are sufficient
but conservative.
To support more complex code patterns or speculation modes, more sophisticated tools like Kasper \cite{johannesmeyer2022kasper}, KleeSpectre \cite{wang2020kleespectre}, and Spectector \cite{guarnieri2020spectector} could potentially be adapted.
These tools usually do not scale well to large code bases and assume simplified speculation models.
These two limitations are circumvented by the design principles we presented.
\\

\subsection{Burst Mode Security Analysis}
\label{sec:burstsecanalysis}
The attacker observations are limited to \sw{shm}. 
Burst mode only allows for straight-line speculation, and control flow is self-contained.
That means that all reachable speculative paths are described by \sw{stl}. 
Under our threat model, $\uarch[burst]{\cdot}{} \vdash \contract[shm]{\cdot}{stl}$.
\\

\noindent\textbf{Extending Hardware Semantic With Static Software Analysis}
We extend hardware semantics to capture a static-software-analysis step.
Let's consider a static analysis tool \hw{tool} that takes any program $p$ as an input and returns \texttt{true} if the program passes the static analysis and \texttt{false} otherwise.
Given a static analysis program \hw{tool}, and a hardware semantic $\uarch[]{\cdot}{}$, for all program $p$, we define
\begin{align*}
    \uarch[]{p}{tool} = \begin{cases}
        \uarch[]{p}{} & \text{if } tool(p) = \texttt{true} \\
        \varnothing & \text{otherwise}
    \end{cases}
\end{align*}
Where $\varnothing$ represents the hardware semantic that only returns empty traces.
This can be easily implemented by adding a condition in the SM that verifies the (precomputed) result of \hw{tool}($p$) before starting the enclave.
An interesting special static analysis program is $\hw{\top}$ which returns true for any program.
\\

\noindent\textbf{Qualifying Our Static Analysis}
We call the static analysis used for Burst mode \hw{sta}.
Our tool assumes the \sw{shm} leakage model and only accepts programs that will not leak more secrets under straight-line speculation \sw{stl} than under sequential execution \sw{seq}.
That is, for a program $p$, \hw{sta} returns \texttt{true} if $p \vdash NI(\contract[shm]{\cdot}{seq} \Rightarrow \contract[shm]{\cdot}{stl})$, and \texttt{false} otherwise.
\\

\noindent
This gives us $\uarch[burst]{p}{sta} \vdash \contract[shm]{\cdot}{seq}$ (see proof in appendix~\ref{app:proofburstmode}).
In summary, \textbf{Burst mode satisfies RMI.}
\\

\noindent\textbf{Secure Composability}
Because our code snippets (speculative) control flows are self-contained, they can be safely reused and composed with code using other execution modes.
Developers can leverage pre-analyzed performance libraries (e.g., \texttt{memcpy}) without additional instrumentation or analysis.
In our design, we show that Burst mode on key code snippets significantly improves overall performance (see Section \ref{sec:evalsharedmemory}).